
\documentclass[11pt,a4paper]{article}

% ==================== PAQUETES ====================
\usepackage[utf8]{inputenc}
\usepackage[T1]{fontenc}
\usepackage[spanish]{babel}
\usepackage[margin=2cm]{geometry}
\usepackage{tabularx}
\usepackage{array}
\usepackage{longtable}
\usepackage{booktabs}
\usepackage{enumitem}
\usepackage{setspace}
\usepackage{titlesec}
\usepackage[table]{xcolor}
\usepackage{graphicx}
\usepackage{float}
\usepackage{hyperref}

% ==================== CONFIGURACIÓN ====================
\definecolor{azulTabla}{RGB}{31,78,121}

\hypersetup{
	colorlinks=true,
	linkcolor=azulTabla,
	citecolor=azulTabla,
	urlcolor=azulTabla
}

\newcolumntype{Y}{>{\raggedright\arraybackslash}X}

\setstretch{1.15}
\setlength{\parindent}{0pt}
\setlength{\parskip}{6pt}

\titleformat{\section}
{\normalfont\Large\bfseries\color{azulTabla}}
{\thesection.}{0.5em}{}

\titleformat{\subsection}
{\normalfont\large\bfseries}
{\thesubsection}{0.5em}{}

% ==================== PLANTILLAS REUTILIZABLES ====================

\newcommand{\datosclave}{
	\begin{tabularx}{\textwidth}{|p{4cm}|Y|}
		\hline
		\textbf{Dato clave} & \textbf{Detalle} \\ \hline
		Equipo & \\ \hline
		Plazo & \\ \hline
		Disponibilidad del cliente & \\ \hline
		Cambios esperados & \\ \hline
		Criticidad & \\ \hline
		Marco normativo o contractual & \\ \hline
	\end{tabularx}
}

\newcommand{\matriz}{
	\begin{tabularx}{\textwidth}{|p{2.6cm}|Y|Y|Y|}
		\hline
		\textbf{Criterio} &
		\textbf{Secuencial (cascada)} &
		\textbf{Iterativo e incremental} &
		\textbf{Ágil} \\ \hline
		K1. Estabilidad y cambios & & & \\[2cm] \hline
		K2. Partes interesadas y disponibilidad & & & \\[2cm] \hline
		K3. Tipo de sistema y criticidad & & & \\[2cm] \hline
		K4. Riesgo principal & & & \\[2cm] \hline
		K5. Documentación y marco normativo o contractual & & & \\[2cm] \hline
	\end{tabularx}
}

\newcommand{\justificacion}{
	\begin{tabularx}{\textwidth}{|p{4.5cm}|Y|}
		\hline
		\textbf{Elemento} & \textbf{Respuesta} \\ \hline
		Modelo recomendado (puro o híbrido) & \\[1.3cm] \hline
		Razón 1 (dato, criterio, consecuencia) & \\[1.3cm] \hline
		Razón 2 & \\[1.3cm] \hline
		Razón 3 & \\[1.3cm] \hline
		Por qué se descartan los otros dos modelos & \\[1.3cm] \hline
		Tensión principal del escenario & \\[1.3cm] \hline
	\end{tabularx}
}

\newcommand{\riesgos}{
	\begin{tabularx}{\textwidth}{|p{1.5cm}|Y|Y|}
		\hline
		\textbf{N.º} & \textbf{Riesgo del modelo elegido} & \textbf{Mitigación} \\ \hline
		1 & & \\[1.5cm] \hline
		2 & & \\[1.5cm] \hline
	\end{tabularx}
}

% ==================== DOCUMENTO ====================

\begin{document}
	
	% ------------------------------------------------
	% CARÁTULA ORIGINAL
	% ------------------------------------------------
	
	\begin{titlepage}
		\centering
		
		{\large\textbf{UNIVERSIDAD TÉCNICA ESTATAL DE QUEVEDO}}\\[0.3cm]
		
		{\normalsize Facultad de Ciencias de la Ingeniería}\\[0.15cm]
		{\normalsize Carrera de Ingeniería de Software}
		
		\vspace{1.5cm}
		
		{\color{azulTabla}\rule{0.85\textwidth}{1.5pt}}
		
		\vspace{1.2cm}
		
		{\large\textcolor{azulTabla}{\textbf{ACTIVIDAD ACADÉMICA}}}\\[0.5cm]
		
		{\Large\textbf{FORM-PC-01}}\\[0.8cm]
		
		{\LARGE\bfseries
			Comparación de modelos\\[0.2cm]
			de proceso de requisitos
		}
		
		\vspace{1.2cm}
		
		{\color{azulTabla}\rule{0.85\textwidth}{1.5pt}}
		
		\vspace{0.8cm}
		
		\begin{minipage}{0.85\textwidth}
			\renewcommand{\arraystretch}{1.7}
			
			\begin{tabular}{@{}>{\bfseries}p{3.2cm}p{9cm}@{}}
				
				Asignatura: &
				Ingeniería de Requerimientos (ISR-401) \\
				
				Código: &
				ISR-401 \\
				
				Docente: &
				Gleiston Cicerón Guerrero Ulloa, PhD \\
				
				Integrantes: &
				HOLGUIN MENDOZA JONATHAN ALEXANDER \\
				
				&
				LAVAYEN BONE CRISTOPHER JAVIER \\
				
				&
				MORAN PILAGUANO FRIXON FERNANDO \\
				
				&
				SALTOS TELLO JOSEPH ORLANDO \\
				
				Paralelo: &
				CUARTO SOFTWARE "B" \\
				
				Fecha de entrega: &
				12/11/2026 \\
				
			\end{tabular}
		\end{minipage}
		
		\vfill
		
		{\color{azulTabla}\rule{0.85\textwidth}{0.8pt}}\\[0.4cm]
		
		{\normalsize\textbf{QUEVEDO -- LOS RÍOS -- ECUADOR}}\\[0.15cm]
		{\normalsize 2026}
		
	\end{titlepage}
	
	% ==================== 1 ====================
	
	\section{Introducción}
	
La elección de un modelo de proceso de requisitos —secuencial, iterativo o ágil— condiciona el éxito de un proyecto de software y depende de factores como su estabilidad, criticidad y la disponibilidad de los interesados. En este trabajo se comparan dichos modelos frente a cuatro escenarios reales, apoyándose en referentes conceptuales y normativos como la norma ISO/IEC/IEEE 29148, SWEBOK y el framework de Klaus Pohl para determinar el enfoque idóneo según cada entorno técnico y contractual.
	% ==================== 2 ====================
	
	\section{Objetivos}
	
	\subsection{Objetivo general}
	
	Comparar los modelos de proceso de requisitos en escenarios reales para aprender a elegir el más adecuado según el tipo de sistema, la estabilidad de los requisitos y la disponibilidad de las partes interesadas.
	
	\subsection{Objetivos específicos}
	
	\begin{itemize}[noitemsep]
	\item Caracterizar las familias de modelos de proceso de requisitos identificando sus fases, supuestos y condiciones de uso.
	\item Evaluar el comportamiento de cada modelo mediante una matriz comparativa contextualizada a partir de criterios técnicos en cuatro escenarios reales.
	\item Justificar la recomendación de un modelo de proceso para cada escenario, identificando sus riesgos específicos y respectivas estrategias de mitigación.
	\item Describir los actores del proceso y ubicar las actividades centrales, transversales y artefactos del framework de Klaus Pohl en un escenario seleccionado.
\end{itemize}
	
	% ==================== 3 ====================
	
	\section{Descripción de los tres modelos de proceso}
	
	\subsection{Modelo secuencial (cascada)}
	
	\begin{tabularx}{\textwidth}{|p{3.5cm}|Y|}
		\hline
		\textbf{Elemento} & \textbf{Descripción} \\ \hline
		Fases &
		Se desarrolla mediante etapas ordenadas: levantamiento y análisis de requisitos, diseño, implementación, pruebas, despliegue y mantenimiento. Cada fase se completa antes de avanzar a la siguiente. \\ \hline
		Supuestos &
		Los requisitos pueden definirse con suficiente claridad desde el inicio y se espera que permanezcan relativamente estables durante el desarrollo. \\ \hline
		Condiciones de uso &
		Es apropiado cuando existe poca incertidumbre, los cambios son limitados y se necesita una planificación y documentación detalladas desde las primeras etapas. \\ \hline
	\end{tabularx}
	
	\subsection{Modelo iterativo e incremental}
	
	\begin{tabularx}{\textwidth}{|p{3.5cm}|Y|}
		\hline
		\textbf{Elemento} & \textbf{Descripción} \\ \hline
		Fases &
		El sistema se construye mediante iteraciones e incrementos. En cada ciclo se realizan actividades de análisis, diseño, implementación y pruebas, incorporando nuevas funcionalidades o mejorando las existentes. \\ \hline
		Supuestos &
		Los requisitos pueden evolucionar y no necesariamente se conocen por completo desde el inicio. Se considera posible entregar partes funcionales del sistema progresivamente. \\ \hline
		Condiciones de uso &
		Es adecuado cuando se necesita reducir la incertidumbre gradualmente, obtener resultados parciales y realizar ajustes entre ciclos de desarrollo. \\ \hline
	\end{tabularx}
	
	\subsection{Modelo ágil}
	
	\begin{tabularx}{\textwidth}{|p{3.5cm}|Y|}
		\hline
		\textbf{Elemento} & \textbf{Descripción} \\ \hline
		Fases &
		El trabajo se organiza en ciclos cortos. Se priorizan requisitos, se planifica una iteración, se desarrolla y prueba un incremento, y se obtiene retroalimentación para ajustar el siguiente ciclo. \\ \hline
		Supuestos &
		Los requisitos pueden cambiar durante el proyecto y la participación frecuente de los interesados permite revisar prioridades y validar las funcionalidades desarrolladas. \\ \hline
		Condiciones de uso &
		Es apropiado cuando existe incertidumbre o alta volatilidad de requisitos, se necesita entregar valor rápidamente y los usuarios pueden participar de manera continua. \\ \hline
	\end{tabularx}
	
	% ==================== 4 ====================
	
	\section{Criterios de comparación}
	
	\begin{tabularx}{\textwidth}{|p{4.5cm}|Y|}
		\hline
		\textbf{Criterio} & \textbf{Qué significa} \\ \hline
		K1. Estabilidad y gestión de cambios &
		Evalúa qué tan estables son los requisitos y la capacidad del modelo para incorporar cambios sin generar costos, retrasos o retrabajo excesivo. \\ \hline
		K2. Partes interesadas y su disponibilidad &
		Considera la participación de clientes, usuarios y responsables del proyecto durante el levantamiento, validación y modificación de los requisitos. \\ \hline
		K3. Tipo de sistema y criticidad &
		Analiza la naturaleza del sistema, sus consecuencias ante fallos y el nivel de control, validación y rigurosidad que requiere. \\ \hline
		K4. Riesgo principal del proyecto &
		Identifica la amenaza más importante para el cumplimiento de los objetivos, como cambios frecuentes, errores, retrasos, sobrecostos o falta de validación. \\ \hline
		K5. Documentación y marco normativo o contractual &
		Valora la necesidad de documentar requisitos, decisiones y pruebas, además de cumplir obligaciones legales, normativas o contractuales. \\ \hline
	\end{tabularx}
	
	\textbf{Escala de valoración:} para las matrices comparativas se utiliza una escala de 1 a 5, donde 1 representa una adaptación muy baja del modelo al criterio y 5 representa una adaptación muy alta. Las puntuaciones corresponden a una valoración cualitativa del equipo según las características de cada escenario.
	
	% ==================== 5 ====================
	
	\section{Análisis de los cuatro escenarios}
	
	\subsection{Escenario 1. Extractora de aceite de palma}
	
	
	\begin{tabularx}{\textwidth}{|p{4cm}|Y|}
	\hline
	\textbf{Dato clave} & \textbf{Detalle} \\ \hline
	Equipo &
	4 personas (2 de software embebido, 1 analista de requisitos y 1 encargado de pruebas). \\ \hline
	Plazo &
	10 meses, con puesta en marcha obligatoria antes de la próxima temporada alta. \\ \hline
	Disponibilidad del cliente &
	Baja a media: el jefe de planta solo puede reunirse dos mañanas por semana; los operadores trabajan en turnos rotativos. \\ \hline
	Cambios esperados &
	Baja volatilidad en funciones de seguridad, alta volatilidad en la interfaz del operador, reportes y la comunicación con sensores aún no comprados. \\ \hline
	Criticidad &
	Muy alta: involucra la seguridad física de los operadores (riesgo de quemaduras y presión residual), prevención de desbordes de tanques y control de pérdidas de producción. \\ \hline
	Marco normativo o contractual &
	Exigencias estrictas de la aseguradora de requisitos de seguridad aprobados y pruebas documentadas obligatorias antes de operar y especificaciones técnicas fijadas por el fabricante de los esterilizadores. \\ \hline
\end{tabularx}
	
	\subsection{Escenario 2. Aplicación móvil de pedidos}
	
	\begin{tabularx}{\textwidth}{|p{4cm}|Y|}
		\hline
		\textbf{Dato clave} & \textbf{Detalle} \\ \hline
		Equipo &
		Pequeño equipo de desarrollo y propietario del negocio como responsable de validar las necesidades. \\ \hline
		Plazo &
		14 semanas para desarrollar y entregar la aplicación móvil. \\ \hline
		Disponibilidad del cliente &
		El propietario está disponible para participar, revisar avances y proporcionar retroalimentación. \\ \hline
		Cambios esperados &
		Alta, debido a que el propietario cambia frecuentemente sus ideas y prioridades sobre las funcionalidades. \\ \hline
		Criticidad &
		Media. La aplicación debe gestionar correctamente pedidos, pagos y entregas, aunque no corresponde a un sistema industrial de seguridad crítica. \\ \hline
		Marco normativo o contractual &
		Se deben considerar las obligaciones aplicables sobre protección de datos personales, privacidad y tratamiento de información relacionada con pedidos y pagos. \\ \hline
	\end{tabularx}
	
	\subsection{Escenario 3. Sistema de trámites municipales}

	El GAD Municipal de San Isidro del Baba, en Los Ríos, es un municipio
	ficticio que atiende cerca de 9\,000 trámites al año. Entre ellos están
	los permisos de funcionamiento, las patentes y los certificados de no
	adeudar. Actualmente, los ciudadanos pueden tener que ir hasta tres
	veces a ventanilla y los técnicos revisan carpetas en papel. El municipio
	contrató un sistema web para hacer estos trámites en línea.

	El sistema se instalará en los servidores del municipio y deberá
	consultar los datos de los predios en una base antigua que carece de
	documentación.

	\begin{tabularx}{\textwidth}{|p{4cm}|Y|}
		\hline
		\textbf{Dato clave} & \textbf{Detalle} \\ \hline
		Equipo & Seis personas de la empresa contratista. \\ \hline
		Plazo & Siete meses. En el mes 2 debe quedar aprobada la
		especificación de requisitos. En el mes 6 se entrega el sistema
		instalado y probado. En el mes 7 se realiza la capacitación y se
		firma el acta de entrega-recepción definitiva. \\ \hline
		Disponibilidad del cliente & El administrador del contrato se reúne
		con el equipo una vez por semana y es el único contacto oficial.
		Los técnicos de ventanilla y la jefa de rentas son los usuarios
		principales y participan cuando él los convoca. \\ \hline
		Cambios esperados & La jefa de rentas ya cambió dos veces los pasos
		para aprobar patentes. El concejo municipal podría aprobar una
		ordenanza en el mes 4 que cambiaría la forma de calcular las tasas.
		\\ \hline
		Criticidad & Alta. El sistema manejará datos personales y cobros
		municipales. Un error puede afectar el dinero público y generar
		reclamos. \\ \hline
		Marco normativo o contractual & El contrato público fija las fechas
		y los pagos de los tres entregables. Hay una multa por cada día de
		retraso. Si un cambio de alcance aumenta el costo, requiere un
		trámite formal que puede tardar semanas. \\ \hline
	\end{tabularx}

	\subsection{Escenario 4. Plataforma para cooperativa de cacao}
	
	\begin{tabularx}{\textwidth}{|p{4cm}|Y|}
		\hline
		\textbf{Dato clave} & \textbf{Detalle} \\ \hline
		Equipo & 4 personas. \\ \hline
		Plazo & 12 meses, con un informe de avance al donante cada semestre. \\ \hline
		Disponibilidad del cliente & Mixta. La gerente atiende casi a diario, pero no decide sola; la directiva se reúne una vez al mes; los 640 socios solo están disponibles los días de acopio y en las asambleas trimestrales. \\ \hline
		Cambios esperados & Bajos en las reglas de liquidación, fijadas en el estatuto. Altos en la trazabilidad y la certificación del cacao, que los compradores modifican cada temporada. \\ \hline
		Criticidad & Alta: maneja el dinero y los créditos de los socios, y sus datos personales. \\ \hline
		Marco normativo o contractual & Estatuto de la cooperativa, control del organismo que supervisa la economía popular y solidaria, condiciones del donante y normativa de protección de datos personales. \\ \hline
	\end{tabularx}
	
	% ==================== 6 ====================
	
	\section{Matrices comparativas}
	
	\subsection{Escenario 1. Extractora de aceite de palma}
	
	\begin{longtable}{|p{2.6cm}|*{3}{>{\raggedright\arraybackslash}p{\dimexpr(\textwidth-2.6cm-8\tabcolsep-5\arrayrulewidth)/3\relax}|}}
	\hline
	\textbf{Criterio} &
	\textbf{Secuencial (cascada)} &
	\textbf{Iterativo e incremental} &
	\textbf{Ágil} \\ \hline
	\endfirsthead
	\hline
	\textbf{Criterio} &
	\textbf{Secuencial (cascada)} &
	\textbf{Iterativo e incremental} &
	\textbf{Ágil} \\ \hline
	\endhead
	K1. Estabilidad y cambios &
	\textbf{4/5.} Las funciones de seguridad y las reglas del fabricante son estables y no cambiarán. Sin embargo, los sensores aún no se compran y la interfaz del operador es incierta, lo que obligaría a aplicar un control estricto de cambios sobre la línea base. &
	\textbf{3/5.} Permite estabilizar primero las funciones críticas de seguridad y dejar para ciclos posteriores la interfaz o los reportes, aunque la falta de definición de los sensores complica la planificación por incrementos fijos. &
	\textbf{2/5.} La alta volatilidad de la interfaz y la integración de sensores no comprados choca con la baja disponibilidad del jefe de planta (dos mañanas por semana), dificultando la priorización continua de un marco puramente ágil. \\ \hline
	K2. Partes interesadas y disponibilidad &
	\textbf{3/5.} Requiere reuniones intensas al inicio para congelar las funciones de seguridad con el jefe de planta y el responsable de seguridad industrial, encajando bien con su disponibilidad limitada de dos mañanas por semana. &
	\textbf{3/5.} Exige revisiones coordinadas al cierre de cada ciclo. La disponibilidad intermitente del jefe de planta y los turnos rotativos de los operadores dificultan validaciones frecuentes sin una planificación formal previa. &
	\textbf{1/5.} Es sumamente difícil de aplicar porque exige una participación constante y fluida del cliente, la cual no existe debido a que el jefe de planta solo atiende dos mañanas semanales y los operadores laboran por turnos. \\ \hline
	K3. Tipo de sistema y criticidad &
	\textbf{5/5.} Sistema embebido industrial de criticidad muy alta (riesgo de quemaduras por presión y desbordes). La rigurosidad del modelo en la especificación formal y la verificación exhaustiva responde perfectamente a la exigencia de la aseguradora. &
	\textbf{4/5.} Permite segregar los componentes críticos de seguridad y desarrollarlos bajo controles estrictos, incorporando los incrementos de la interfaz de forma controlada y validando cada parte antes de la puesta en marcha. &
	\textbf{2/5.} Evaluar componentes críticos de seguridad industrial (presión de vapor y bloqueos de compuertas) mediante ciclos cortos informales incrementa el riesgo de omisiones normativas exigidas por la aseguradora. \\ \hline
	K4. Riesgo principal &
	\textbf{3/5.} El mayor riesgo es descubrir tarde problemas en la interfaz o incompatibilidades con los sensores no comprados, lo que implicaría costos altos de rediseño tras aprobar la línea base inicial. &
	\textbf{4/5.} Mitiga el riesgo de integración por fases, permitiendo probar partes funcionales del controlador de forma progresiva antes de la temporada alta de los 10 meses de plazo. &
	\textbf{3/5.} El riesgo principal es descuidar los requisitos no funcionales críticos (memoria limitada, red local sin internet) y la documentación obligatoria exigida por la póliza de seguros de la extractora. \\ \hline
	K5. Documentación y marco normativo o contractual &
	\textbf{5/5.} Se alinea totalmente con las exigencias de la aseguradora, que demanda un documento de requisitos de seguridad rigurosamente aprobado y evidencias documentadas de pruebas de cada función de protección antes de la puesta en marcha. &
	\textbf{4/5.} Permite generar la documentación y las evidencias de prueba de forma progresiva conforme se cierran los incrementos de seguridad y control de tanques exigidos por la normativa de la planta. &
	\textbf{2/5.} Tiende a minimizar la documentación formal, lo cual choca frontalmente con la exigencia jurídica y contractual de la aseguradora de contar con manuales y especificaciones técnicas aprobadas de antemano. \\ \hline
\end{longtable}
	
	\subsection{Escenario 2. Aplicación móvil de pedidos}
	
	\noindent\begin{minipage}{\textwidth}
	\begin{center}
		\textbf{Matriz comparativa del escenario 2: aplicación móvil de pedidos}
	\end{center}
	
	\begin{tabularx}{\textwidth}{|p{2.6cm}|Y|Y|Y|}
		\hline
		\textbf{Criterio} &
		\textbf{Secuencial (cascada)} &
		\textbf{Iterativo e incremental} &
		\textbf{Ágil} \\ \hline
		
		K1. Estabilidad y cambios &
		\textbf{1/5}. Presenta baja adaptación porque los cambios frecuentes pueden provocar retrabajo y modificaciones importantes en etapas ya terminadas. &
		\textbf{4/5}. Permite incorporar modificaciones entre iteraciones y ajustar progresivamente las funcionalidades. &
		\textbf{5/5}. Facilita adaptar prioridades y requisitos mediante ciclos cortos y retroalimentación continua. \\ \hline
		
		K2. Partes interesadas y disponibilidad &
		\textbf{2/5}. La participación se concentra principalmente en las etapas iniciales y en las revisiones establecidas. &
		\textbf{4/5}. Permite consultar al propietario al finalizar cada incremento y considerar sus observaciones. &
		\textbf{5/5}. Aprovecha la disponibilidad del propietario para validar continuamente las funcionalidades y prioridades. \\ \hline
		
		K3. Tipo de sistema y criticidad &
		\textbf{3/5}. Puede desarrollar una aplicación de este tipo, pero su rigidez dificulta responder a cambios funcionales. &
		\textbf{4/5}. Permite construir y probar progresivamente los módulos de pedidos, pagos y entregas. &
		\textbf{5/5}. Facilita entregar primero las funciones esenciales y mejorar la aplicación según las necesidades del negocio. \\ \hline
		
		K4. Riesgo principal &
		\textbf{2/5}. Existe riesgo de entregar una aplicación que no responda a las expectativas cambiantes del propietario. &
		\textbf{4/5}. Reduce el riesgo mediante entregas parciales, aunque requiere controlar el alcance de cada incremento. &
		\textbf{5/5}. Permite detectar diferencias entre lo solicitado y lo desarrollado mediante validaciones frecuentes. \\ \hline
		
		K5. Documentación y marco normativo o contractual &
		\textbf{4/5}. Facilita establecer documentación y requisitos formales desde el inicio. &
		\textbf{4/5}. Permite mantener documentación actualizada durante los incrementos del sistema. &
		\textbf{3/5}. Prioriza el software funcional, por lo que se debe establecer expresamente la documentación y los controles de privacidad y seguridad necesarios. \\ \hline
	\end{tabularx}
	\end{minipage}\par
	
	\clearpage
	\subsection{Escenario 3. Sistema de trámites municipales}

	\begin{longtable}{|p{2.6cm}|*{3}{>{\raggedright\arraybackslash}p{\dimexpr(\textwidth-2.6cm-8\tabcolsep-5\arrayrulewidth)/3\relax}|}}
		\hline
		\textbf{Criterio} &
		\textbf{Secuencial (cascada)} &
		\textbf{Iterativo e incremental} &
		\textbf{Ágil} \\ \hline
		\endfirsthead
		\hline
		\textbf{Criterio} &
		\textbf{Secuencial (cascada)} &
		\textbf{Iterativo e incremental} &
		\textbf{Ágil} \\ \hline
		\endhead
		K1. Estabilidad y cambios &
		\textbf{2/5.} La aprobación de patentes ya cambió dos veces y las tasas podrían
		cambiar en el mes 4. Para entonces los requisitos ya estarían aprobados, así que
		tocaría revisar lo acordado y rehacer parte del trabajo. &
		\textbf{4/5.} Se puede revisar el flujo de patentes en cada ciclo y ajustar
		el cálculo si la ordenanza se aprueba en el mes 4. El alcance aprobado
		sigue sujeto al contrato. &
		\textbf{4/5.} Los ciclos cortos permiten probar los cambios de patentes
		y dar prioridad a las nuevas tasas si se aprueban. Los cambios que
		aumenten el costo requieren un trámite formal. \\ \hline
		K2. Partes interesadas y disponibilidad &
		\textbf{3/5.} La reunión semanal con el administrador permite programar las
		revisiones. Como los técnicos y la jefa de rentas solo participan cuando él los
		convoca, pueden quedar dudas sobre los trámites antes de aprobar los requisitos
		en el mes 2. &
		\textbf{4/5.} Las revisiones de cada avance se pueden coordinar en la
		reunión semanal con el administrador. Él puede convocar a la jefa
		de rentas para revisar las patentes. &
		\textbf{2/5.} El administrador es el único contacto oficial y se reúne
		una vez por semana. Si los técnicos solo participan por convocatoria,
		algunas dudas pueden quedar abiertas entre ciclos. \\ \hline
		K3. Tipo de sistema y criticidad &
		\textbf{4/5.} El sistema web maneja datos personales y cobra tasas. La fase de
		pruebas permite revisar los cálculos y quién puede consultar los datos antes de
		entregar el sistema. &
		\textbf{5/5.} Cada avance permite probar los cobros y el acceso a datos
		personales antes de sumar más trámites. Así se pueden detectar errores
		que afectarían al dinero público. &
		\textbf{4/5.} Las pruebas en cada ciclo ayudan a detectar errores en
		los cobros y en el acceso a datos personales. El equipo debe guardar
		los resultados para la entrega del mes 6. \\ \hline
		K4. Riesgo principal &
		\textbf{2/5.} La base de datos de los predios es antigua y no está documentada.
		Si el equipo espera hasta las pruebas finales para comprobar que puede leerla,
		un problema de conexión o de datos podría retrasar la entrega del mes 6 y
		generar multas. &
		\textbf{5/5.} Un primer avance puede probar la lectura de la base
		antigua desde los servidores del municipio. Así se detectan problemas
		con tiempo para atenderlos antes de la entrega del mes 6. &
		\textbf{4/5.} Un ciclo corto al inicio permite probar la lectura de
		la base sin documentación. El equipo puede ajustar sus prioridades
		si descubre problemas antes de la entrega del mes 6. \\ \hline
		K5. Documentación y marco normativo o contractual &
		\textbf{5/5.} El trabajo por fases encaja con los entregables y pagos fijados
		para los meses 2, 6 y 7. Los requisitos aprobados en el mes 2 sirven como base
		para documentar los cambios de alcance que aumenten el costo y requieran un
		trámite formal. &
		\textbf{4/5.} Se pueden revisar los avances entre las entregas de los
		meses 2, 6 y 7. Cada ciclo debe mantener al día los requisitos
		aprobados y los resultados de las pruebas para la entrega del mes 6. &
		\textbf{2/5.} El contrato exige aprobar los requisitos en el mes 2.
		Cambiar prioridades en un ciclo resulta útil, pero si el cambio
		aumenta el costo, el equipo debe esperar el trámite formal. \\ \hline
	\end{longtable}
	
	\clearpage
	\subsection{Escenario 4. Plataforma para cooperativa de cacao}
	
	\begin{center}
		\textbf{Matriz comparativa del escenario 4: plataforma para cooperativa de cacao}
	\end{center}
	
	\small
	\begin{tabularx}{\textwidth}{|p{2.2cm}|Y|Y|Y|}
		\hline
		\textbf{Criterio} &
		\textbf{Secuencial (cascada)} &
		\textbf{Iterativo e incremental} &
		\textbf{Ágil} \\ \hline
		
		K1. Estabilidad y cambios &
		\textbf{2/5}. Las liquidaciones salen del estatuto y casi no cambian, así que se pueden fijar al inicio. Las exigencias de trazabilidad y certificación que los compradores modifican cada temporada sí cambian, y cada ajuste pasaría por control formal sobre la línea base. &
		\textbf{4/5}. Se acuerda una base con liquidaciones, entregas y saldos. Lo que pidan los compradores en cada temporada entra como trabajo de un ciclo posterior, sin tocar lo ya entregado. &
		\textbf{4/5}. Reordenar una lista priorizada en cada ciclo ayuda con la trazabilidad, pero las liquidaciones, que el estatuto mantiene fijas, no tienen cambios que absorber. \\ \hline
		
		K2. Partes interesadas y disponibilidad &
		\textbf{3/5}. La gerente (casi a diario) y la directiva (una vez al mes) alcanzan para aprobar cada fase. Los socios, que hoy no entienden cómo se calcula su pago, solo verían el sistema ya terminado. &
		\textbf{4/5}. Cada ciclo puede cerrar con una revisión en la reunión mensual de la directiva y con una prueba en un día de acopio, que son los momentos en que hay gente disponible. &
		\textbf{2/5}. Requiere a alguien con autoridad para priorizar en cada ciclo. La gerente no decide sola y la directiva solo se reúne cada mes, de modo que las prioridades esperarían semanas. \\ \hline
		
		K3. Tipo de sistema y criticidad &
		\textbf{3/5}. Un error en el dinero o en el crédito de un socio es grave, y detallar todo antes de construir reduce ese riesgo. El registro de entregas sin conexión en el centro de Mocache, en cambio, solo se comprobaría al final. &
		\textbf{4/5}. Lo crítico (entregas, liquidaciones, saldos) se construye y verifica primero. El registro sin conexión con su sincronización puede probarse en el centro de acopio desde los primeros ciclos. &
		\textbf{2/5}. Los ciclos cortos con poca documentación son riesgosos cuando se maneja dinero de los socios. Los requisitos de seguridad, de datos personales y de uso sin conexión suelen quedar sin escribir. \\ \hline
		
		K4. Riesgo principal &
		\textbf{2/5}. Que lo especificado al comienzo no coincida con lo que pidan los compradores en la próxima temporada, o que resulte difícil de usar para socios con poca experiencia en tecnología. &
		\textbf{4/5}. Repartir mal el alcance entre incrementos y tropezar al unir el módulo sin conexión con el resto del sistema. &
		\textbf{2/5}. Quedarse sin un cliente que decida y sin manera de consultar seguido a los socios, que viven en recintos lejanos. \\ \hline
		
		K5. Documentación y marco normativo o contractual &
		\textbf{3/5}. El estatuto y el control del organismo supervisor se atienden bien con una especificación aprobada. El donante pide funciones terminadas cada semestre, y con una sola entrega final no habría nada que mostrarle a mitad del proyecto. &
		\textbf{4/5}. La especificación crece por versiones y cada semestre sale un incremento terminado con su documentación, que es lo que puede reportarse al donante. &
		\textbf{3/5}. Las historias de usuario y los criterios de aceptación bastan para el equipo, pero quizá no alcancen para rendir cuentas del dinero de los socios ni para demostrar el cumplimiento de la normativa de datos personales. \\ \hline
	\end{tabularx}
	\normalsize
	
	% ==================== 7 ====================
	
	\section{Justificación de los modelos seleccionados}
	
	\subsection{Escenario 1. Extractora de aceite de palma}
	
	\begin{tabularx}{\textwidth}{|p{4.5cm}|Y|}
	\hline
	\textbf{Elemento} & \textbf{Respuesta} \\ \hline
	
	Modelo recomendado (puro o híbrido) &
	Modelo híbrido: Enfoque secuencial (cascada) estricto para los componentes críticos de seguridad y el control de esterilizadores (regulados por el fabricante y exigidos por la aseguradora), combinado con un enfoque iterativo e incremental para la interfaz del operador, reportes y la integración de los sensores pendientes de compra. \\ \hline
	
	Razón 1 (dato, criterio, consecuencia) &
	Las reglas de seguridad de los esterilizadores están definidas por el fabricante y la aseguradora exige requisitos de seguridad aprobados de antemano (dato); esto corresponde a la criticidad y marco normativo (K3, K5), lo que hace indispensable una especificación inicial formal y congelada para evitar riesgos físicos graves en la planta. \\ \hline
	
	Razón 2 &
	El jefe de planta solo puede reunirse dos mañanas por semana y los operadores trabajan en turnos rotativos (dato); esto afecta directamente la disponibilidad de las partes interesadas (K2), imposibilitando un modelo ágil que exige validaciones y retroalimentación continua. \\ \hline
	
	Razón 3 &
	El proveedor de los sensores y válvulas no ha sido elegido y se comprará en el tercer mes, además de existir alta volatilidad en los reportes (dato); esto impacta la estabilidad y gestión de cambios (K1), justificando bloques iterativos posteriores para integrar dichos dispositivos sin alterar la base de seguridad. \\ \hline
	
	Por qué se descartan los otros dos modelos &
	El modelo ágil puro se descarta por completo debido a la baja disponibilidad del cliente, la alta criticidad física de la planta y la exigencia estricta de documentación formal de la aseguradora. El modelo secuencial puro se descarta porque es incapaz de absorber la incertidumbre de los sensores que aún no se compran ni los cambios en la interfaz de pantalla del operador. \\ \hline
	
	Tensión principal del escenario &
	La tensión principal radica en que el sistema combina una parte industrial de seguridad crítica muy estricta (que exige cascada y documentación de aseguradora) con una parte de gestión incierta (sensores no comprados y pantallas por definir, que exigirían flexibilidad). Se resuelve mediante el enfoque híbrido propuesto. \\ \hline
\end{tabularx}
	
	\subsection{Escenario 2. Aplicación móvil de pedidos}
	
	\begin{tabularx}{\textwidth}{|p{4.5cm}|Y|}
		\hline
		\textbf{Elemento} & \textbf{Respuesta} \\ \hline
		
		Modelo recomendado (puro o híbrido) &
		Modelo ágil, con entregas incrementales y validación frecuente del propietario. \\ \hline
		
		Razón 1 (dato, criterio, consecuencia) &
		La volatilidad de los requisitos es alta porque el propietario cambia frecuentemente sus ideas. El modelo ágil permite reorganizar las prioridades y adaptar las funcionalidades durante el desarrollo. \\ \hline
		
		Razón 2 &
		El propietario está disponible para participar y revisar los avances. Esto facilita obtener retroalimentación frecuente y comprobar si las funcionalidades responden a las necesidades reales del negocio. \\ \hline
		
		Razón 3 &
		El plazo de 14 semanas y el presupuesto reducido requieren priorizar las funcionalidades esenciales. El desarrollo ágil permite construir un producto mínimo viable (MVP) y entregar valor progresivamente. \\ \hline
		
		Por qué se descartan los otros dos modelos &
		El modelo secuencial se descarta porque los cambios frecuentes pueden generar retrabajo, retrasos y costos adicionales. El modelo iterativo e incremental es una alternativa viable, pero el modelo ágil ofrece mayor flexibilidad para reorganizar prioridades y recibir retroalimentación continua. \\ \hline
		
		Tensión principal del escenario &
		La principal tensión es equilibrar la flexibilidad para aceptar cambios con el presupuesto limitado y el plazo de 14 semanas. Para controlarla, se debe priorizar un MVP y evaluar el impacto de cada cambio antes de incorporarlo. \\ \hline
	\end{tabularx}
	
	\subsection{Escenario 3. Sistema de trámites municipales}
	
	\begin{tabularx}{\textwidth}{|p{4.5cm}|Y|}
		\hline
		\textbf{Elemento} & \textbf{Respuesta} \\ \hline
		Modelo recomendado &
		\textbf{Iterativo e incremental.} El sistema se desarrolla y se prueba por
		partes, con revisiones al final de cada ciclo. Los avances se organizan
		alrededor de las entregas del contrato: requisitos aprobados en el mes 2,
		sistema instalado y probado en el mes 6, y capacitación con el acta de entrega
		en el mes 7. \\ \hline

		Razón 1. Cambios en los requisitos (K1) &
		La aprobación de patentes ya cambió dos veces y las tasas podrían cambiar en
		el mes 4. Revisar los trámites por partes permite ajustar el paso o el cálculo
		afectado. Si el cambio de alcance aumenta el costo, se tramita su aprobación
		antes de incorporarlo. \\ \hline

		Razón 2. Disponibilidad de los usuarios (K2) &
		El administrador se reúne con el equipo una vez por semana. En esas reuniones
		se pueden acordar las revisiones de cada avance y pedir que convoque a la jefa
		de rentas o a los técnicos cuando haga falta revisar un trámite. \\ \hline

		Razón 3. Riesgo de la base de datos (K4) &
		La base de datos de los predios es antigua y carece de documentación. Uno de
		los primeros avances puede comprobar si el sistema logra leerla desde los
		servidores municipales. Detectar un problema en esa etapa deja más margen para
		resolverlo antes de la entrega del mes 6. \\ \hline

		Por qué se descartan los otros modelos &
		\textbf{Secuencial.} Sus fases encajan con las fechas del contrato, pero una
		ordenanza aprobada en el mes 4 podría obligar a revisar los requisitos
		acordados en el mes 2 y a rehacer parte del sistema. \par \textbf{Ágil.} Los
		ciclos cortos ayudan a ajustar prioridades, pero los usuarios participan solo
		cuando los convoca el administrador. Esa disponibilidad limita la revisión
		frecuente y el contrato exige un trámite formal si un cambio de alcance
		aumenta el costo. \\ \hline

		Tensión principal del escenario &
		Los pagos, las fechas fijas y las multas por retraso favorecen un trabajo por
		fases. Al mismo tiempo, los pasos de aprobación y las tasas pueden cambiar
		después de aprobar los requisitos. La propuesta mantiene las fechas del
		contrato y usa los ciclos para probar cada parte y revisar los cambios que
		surjan. \\ \hline
	\end{tabularx}
	
	\subsection{Escenario 4. Plataforma para cooperativa de cacao}
	
	\small
	\begin{tabularx}{\textwidth}{|p{4.5cm}|Y|}
		\hline
		\textbf{Elemento} & \textbf{Respuesta} \\ \hline
		
		Modelo recomendado (puro o híbrido) &
		Híbrido con base iterativa e incremental. El núcleo estable (entregas, liquidaciones, saldos de crédito, avisos) se especifica y la directiva lo aprueba en los primeros meses como línea base. La trazabilidad y certificación de los compradores y el registro sin conexión se construyen en iteraciones, con una revisión al terminar cada una. \\ \hline
		
		Razón 1 (dato, criterio, consecuencia) &
		Dato: las reglas de liquidación constan en el estatuto y no cambian (K1). Consecuencia: pueden escribirse y aprobarse al inicio sin temor a que queden obsoletas, y solo el módulo de trazabilidad necesita ciclos. \\ \hline
		
		Razón 2 &
		Dato: el donante exige un informe cada seis meses con funciones terminadas (K5). Consecuencia: el plan debe producir incrementos verificables en cada semestre; una única entrega final dejaría los informes sin respaldo. \\ \hline
		
		Razón 3 &
		Dato: la gerente atiende casi a diario pero no decide sola, la directiva sesiona una vez al mes y los socios solo aparecen en acopio y asambleas (K2). Consecuencia: no hay quien priorice cada semana, pero sí hay momentos fijos de revisión a los que atar el cierre de cada ciclo. \\ \hline
		
		Por qué se descartan los otros dos modelos &
		Cascada pura: los compradores cambian la trazabilidad cada temporada y una sola entrega final no permite informar al donante a los seis meses. Ágil puro: falta un representante con autoridad para priorizar, y la documentación mínima se queda corta en un sistema que maneja dinero y datos personales. \\ \hline
		
		Tensión principal del escenario &
		La gerente, disponible casi todos los días, y los cambios frecuentes de la trazabilidad empujan hacia un modelo ágil. Frenan ese camino la directiva mensual, la criticidad alta y las reglas fijas del estatuto. \\ \hline
	\end{tabularx}
	\normalsize
	
	% ==================== 8 ====================
	
	\section{Riesgos y mitigaciones}
	
	\subsection{Escenario 1. Extractora de aceite de palma}
	
	\begin{tabularx}{\textwidth}{|p{1.5cm}|Y|Y|}
	\hline
	\textbf{N.º} & \textbf{Riesgo del modelo elegido} & \textbf{Mitigación} \\ \hline
	
	1 &
	\textbf{Incompatibilidad técnica con los nuevos sensores:} Dado que el proveedor de válvulas y sensores se definirá en el tercer mes, las decisiones iniciales de diseño en cascada podrían entrar en conflicto con los protocolos del hardware adquirido posteriormente, retrasando el plazo límite de los 10 meses. &
	Definir una capa de abstracción de hardware (firmware modular) durante el análisis inicial que independice la lógica de control principal de las especificaciones físicas específicas del sensor que se elija en el tercer mes. \\ \hline
	
	2 &
	\textbf{Incumplimiento de las exigencias documentales de la aseguradora:} Si la documentación formal de los requisitos de seguridad y las evidencias de pruebas no se estructuran rigurosamente de acuerdo con las políticas de la póliza antes de la puesta en marcha, la aseguradora podría invalidar el funcionamiento del sistema. &
	Establecer revisiones formales obligatorias de trazabilidad y actas de validación de seguridad con el responsable industrial y la aseguradora al concluir la fase de especificación y en cada incremento crítico antes de implementarlo en la red local. \\ \hline
\end{tabularx}
	
	\subsection{Escenario 2. Aplicación móvil de pedidos}
	
	\begin{tabularx}{\textwidth}{|p{1.5cm}|Y|Y|}
		\hline
		\textbf{N.º} &
		\textbf{Riesgo del modelo elegido} &
		\textbf{Mitigación} \\ \hline
		
		1 &
		\textbf{Crecimiento descontrolado del alcance:} los cambios frecuentes y la incorporación constante de nuevas funcionalidades pueden provocar que el proyecto supere el presupuesto y el plazo de 14 semanas. &
		Definir un producto mínimo viable, mantener una lista priorizada de requisitos y evaluar el impacto de cada cambio en el costo, el tiempo y el alcance antes de aprobarlo. \\ \hline
		
		2 &
		\textbf{Errores en pedidos, pagos y entregas:} Las entregas rápidas y las modificaciones frecuentes pueden ocasionar fallos funcionales o problemas en el tratamiento de información personal. &
		Realizar pruebas funcionales en cada iteración, establecer criterios de aceptación, validar los procesos de pedidos y pagos, y aplicar controles de privacidad y seguridad desde el diseño. \\ \hline
	\end{tabularx}
	
	\subsection{Escenario 2. Aplicación móvil de pedidos}
	
	\begin{tabularx}{\textwidth}{|p{1.5cm}|Y|Y|}
		\hline
		\textbf{N.º} &
		\textbf{Riesgo del modelo elegido} &
		\textbf{Mitigación} \\ \hline
		
		1 &
		\textbf{Crecimiento descontrolado del alcance:} los cambios frecuentes y la incorporación constante de nuevas funcionalidades pueden provocar que el proyecto supere el presupuesto y el plazo de 14 semanas. &
		Definir un producto mínimo viable, mantener una lista priorizada de requisitos y evaluar el impacto de cada cambio en el costo, el tiempo y el alcance antes de aprobarlo. \\ \hline
		
		2 &
		\textbf{Errores en pedidos, pagos y entregas:} Las entregas rápidas y las modificaciones frecuentes pueden ocasionar fallos funcionales o problemas en el tratamiento de información personal. &
		Realizar pruebas funcionales en cada iteración, establecer criterios de aceptación, validar los procesos de pedidos y pagos, y aplicar controles de privacidad y seguridad desde el diseño. \\ \hline
	\end{tabularx}
	
	\subsection{Escenario 3. Sistema de trámites municipales}
	
	\riesgos
	
	\subsection{Escenario 4. Plataforma para cooperativa de cacao}
	
	\small
	\begin{tabularx}{\textwidth}{|p{1.5cm}|Y|Y|}
		\hline
		\textbf{N.º} &
		\textbf{Riesgo del modelo elegido} &
		\textbf{Mitigación} \\ \hline
		
		1 &
		Los incrementos quedan mal repartidos y el módulo sin conexión no encaja con el núcleo del sistema. &
		Fijar en la primera iteración qué entra en cada incremento y con qué criterios se acepta, y probar la sincronización en cada ciclo. \\ \hline
		
		2 &
		Las decisiones se retrasan: la directiva solo se reúne cada mes y la gerente no decide sola. &
		Hacer coincidir el cierre de cada ciclo con la reunión mensual y pactar con la directiva qué cambios menores puede aprobar la gerente. \\ \hline
		
		3 &
		Los socios, con poca experiencia tecnológica, no validan las funciones a tiempo. &
		Llevar prototipos simples a los días de acopio y recoger su opinión antes de cada entrega semestral. \\ \hline
	\end{tabularx}
	\normalsize
	
	% ==================== 9 ====================
	
	\section{Aplicación del framework de Pohl
		(escenario elegido: \underline{\hspace{3cm}})}
	
	\subsection{Actores del proceso}
	
	\begin{tabularx}{\textwidth}{|Y|Y|Y|}
		\hline
		\textbf{Actor del escenario} &
		\textbf{Rol en el proceso de requisitos} &
		\textbf{Cuándo y cómo participa} \\ \hline
		& & \\[1cm] \hline
		& & \\[1cm] \hline
		& & \\[1cm] \hline
		& & \\[1cm] \hline
		& & \\[1cm] \hline
	\end{tabularx}
	
	\subsection{Actividades y artefactos por fase o iteración}
	
	\begin{tabularx}{\textwidth}{|p{2.4cm}|Y|Y|Y|}
		\hline
		\textbf{Fase o iteración} &
		\textbf{Actividades centrales} &
		\textbf{Actividades transversales} &
		\textbf{Artefactos producidos} \\ \hline
		& & & \\[2cm] \hline
		& & & \\[2cm] \hline
		& & & \\[2cm] \hline
	\end{tabularx}
	
	% ==================== 10 ====================
	
	\section{Conclusiones}
	
	% Escribir las conclusiones.
	
\end{document}
